\documentclass[12pt,a4paper]{ctexbook}
\usepackage{geometry}
\geometry{margin=2.2cm}
\usepackage{setspace}
\setstretch{1.35}
\usepackage{hyperref}
\title{曹操诗集}
\author{古典文献整理}
\date{}
\begin{document}
\maketitle
\tableofcontents
\newpage
\section{度关山}
天地间，人为贵。\par
立君牧民，为之轨则。\par
车辙马迹，经纬四极。\par
黜陟幽明，黎庶繁息。\par
於铄贤圣，总统邦域。\par
封建五爵，井田刑狱。\par
有燔丹书，无普赦赎。\par
皋陶甫侯，何有失职？\par
嗟哉后世，改制易律。\par
劳民为君，役赋其力。\par
舜漆食器，畔者十国，\par
不及唐尧，采椽不斫。\par
世叹伯夷，欲以厉俗。\par
侈恶之大，俭为共德。\par
许由推让，岂有讼曲？\par
兼爱尚同，疏者为戚。\par

\section{短歌行 其一}
对酒当歌，人生几何！譬如朝露，去日苦多。\par
慨当以慷，忧思难忘。何以解忧？唯有杜康。\par
青青子衿，悠悠我心。但为君故，沉吟至今。\par
呦呦鹿鸣，食野之苹。我有嘉宾，鼓瑟吹笙。\par
明明如月，何时可掇？忧从中来，不可断绝。(明明 一作：佼佼)\par
越陌度阡，枉用相存。契阔谈讌，心念旧恩。(谈讌 一作：谈宴)\par
月明星稀，乌鹊南飞。绕树三匝，何枝可依？\par
山不厌高，海不厌深。周公吐哺，天下归心。(海 一作：水)\par

\section{短歌行 其二}
周西伯昌，怀此圣德。\par
三分天下，而有其二。\par
修奉贡献，臣节不隆。\par
崇侯谗之，是以拘系。\par
后见赦原，赐之斧钺，得使征伐。\par
为仲尼所称，达及德行，\par
犹奉事殷，论叙其美。\par
齐桓之功，为霸之首。\par
九合诸侯，一匡天下。\par
一匡天下，不以兵车。\par
正而不谲，其德传称。\par
孔子所叹，并称夷吾，民受其恩。\par
赐与庙胙，命无下拜。\par
小白不敢尔，天威在颜咫尺。\par
晋文亦霸，躬奉天王。\par
受赐圭瓒，秬鬯彤弓，\par
卢弓矢千，虎贲三百人。\par
威服诸侯，师之所尊。\par
八方闻之，名亚齐桓。\par
河阳之会，诈称周王，是其名纷葩。\par

\section{对酒}
对酒歌，太平时，吏不呼门。\par
王者贤且明，宰相股肱皆忠良。\par
咸礼让，民无所争讼。\par
三年耕有九年储，仓谷满盈。\par
斑白不负载。\par
雨泽如此，百谷用成。\par
却走马，以粪其土田。\par
爵公侯伯子男，咸爱其民，以黜陟幽明。\par
子养有若父与兄。\par
犯礼法，轻重随其刑。\par
路无拾遗之私。\par
囹圄空虚，冬节不断。\par
人耄耋，皆得以寿终。\par
恩德广及草木昆虫。\par

\section{观沧海}
东临碣石，以观沧海。\par
水何澹澹，山岛竦峙。\par
树木丛生，百草丰茂。\par
秋风萧瑟，洪波涌起。\par
日月之行，若出其中；\par
星汉灿烂，若出其里。\par
幸甚至哉，歌以咏志。\par

\section{龟虽寿}
神龟虽寿，犹有竟时。\par
腾蛇乘雾，终为土灰。(腾 一作：螣)\par
老骥伏枥，志在千里。\par
烈士暮年，壮心不已。\par
盈缩之期，不但在天；\par
养怡之福，可得永年。\par
幸甚至哉，歌以咏志。\par

\section{蒿里行}
关东有义士，兴兵讨群凶。\par
初期会盟津，乃心在咸阳。\par
军合力不齐，踌躇而雁行。\par
势利使人争，嗣还自相戕。\par
淮南弟称号，刻玺於北方。\par
铠甲生虮虱，万姓以死亡。\par
白骨露於野，千里无鸡鸣。\par
生民百遗一，念之断人肠。\par

\section{精列}
厥初生，\par
造划之陶物，莫不有终期。\par
莫不有终期。\par
圣贤不能免，何为怀此忧？\par
愿螭龙之驾，思想昆仑居。\par
思想昆仑居。\par
见期于迂怪，志意在蓬莱。\par
志意在蓬莱。\par
周礼圣徂落，会稽以坟丘。\par
会稽以坟丘。\par
陶陶谁能度？君子以弗忧。\par
年之暮奈何，时过时来微。\par

\section{苦寒行}
北上太行山，艰哉何巍巍！\par
羊肠坂诘屈，车轮为之摧。\par
树木何萧瑟！北风声正悲。\par
熊罴对我蹲，虎豹夹路啼。\par
溪谷少人民，雪落何霏霏！\par
延颈长叹息，远行多所怀。\par
我心何怫郁？思欲一东归。\par
水深桥梁绝，中路正徘徊。\par
迷惑失故路，薄暮无宿栖。\par
行行日已远，人马同时饥。\par
担囊行取薪，斧冰持作糜。\par
悲彼东山诗，悠悠使我哀。\par

\section{气出唱 其一}
华阴山，自以为大。\par
高百丈，浮云为之盖。\par
仙人欲来，出随风，列之雨。\par
吹我洞箫，鼓瑟琴，何訚訚！\par
酒与歌戏，今日相乐诚为乐。\par
玉女起，起舞移数时。\par
鼓吹一何嘈嘈。\par
从西北来时，仙道多驾烟，\par
乘云驾龙，郁何务务。\par
遨游八极，乃到昆仑之山，\par
西王母侧，神仙金止玉亭。\par
来者为谁？赤松王乔，乃德旋之门。\par
乐共饮食到黄昏。\par
多驾合坐，万岁长，宜子孙。\par

\section{气出唱 其二}
驾六龙，乘风而行。\par
行四海，路下之八邦。\par
历登高山临溪谷，乘云而行。\par
行四海外，东到泰山。\par
仙人玉女，下来翱游。\par
骖驾六龙饮玉浆。\par
河水尽，不东流。\par
解愁腹，饮玉浆。\par
奉持行，东到蓬莱山，上至天之门。\par
玉阙下，引见得入，\par
赤松相对，四面顾望，视正焜煌。\par
开玉心正兴，其气百道至。\par
传告无穷闭其口，但当爱气寿万年。\par
东到海，与天连。\par
神仙之道，出窈入冥，常当专之。\par
心恬澹，无所愒。\par
欲闭门坐自守，天与期气。\par
愿得神之人，乘驾云车，\par
骖驾白鹿，上到天之门，来赐神之药。\par
跪受之，敬神齐。\par
当如此，道自来。\par

\section{气出唱 其三}
游君山，甚为真。\par
崔嵬砟硌，尔自为神。\par
乃到王母台，金阶玉为堂，芝草生殿旁。\par
东西厢，客满堂。\par
主人当行觞，坐者长寿遽何央。\par
长乐甫始宜孙子。\par
常愿主人增年，与天相守。\par

\section{却东西门行}
鸿雁出塞北，乃在无人乡。\par
举翅万馀里，行止自成行。\par
冬节食南稻，春日复北翔。\par
田中有转蓬，随风远飘扬。\par
长与故根绝，万岁不相当。\par
奈何此征夫，安得驱四方！\par
戎马不解鞍，铠甲不离傍。\par
冉冉老将至，何时返故乡？\par
神龙藏深泉，猛兽步高冈。\par
狐死归首丘，故乡安可忘！\par

\section{塘上行}
蒲生我池中，其叶何离离。\par
傍能行仁义，莫若妾自知。\par
众口铄黄金，使君生别离。\par
念君去我时，独愁常苦悲。\par
想见君颜色，感结伤心脾。\par
念君常苦悲，夜夜不能寐。\par
莫以豪贤故，弃捐素所爱？\par
莫以鱼肉贱，弃捐葱与薤？\par
莫以麻枲贱，弃捐菅与蒯？\par
出亦复何苦，入亦复何愁。\par
边地多悲风，树木何修修！\par
从君致独乐，延年寿千秋。\par

\section{薤露行}
惟汉廿二世，所任诚不良。\par
沐猴而冠带，知小而谋疆。\par
犹豫不敢断，因狩执君王。\par
白虹为贯日，己亦先受殃。\par
贼臣持国柄，杀主灭宇京。\par
荡覆帝基业，宗庙以燔丧。\par
播越西迁移，号泣而且行。\par
瞻彼洛城郭，微子为哀伤。\par

\section{陌上桑}
驾虹霓，乘赤云，登彼九疑历玉门。\par
济天汉，至昆仑，见西王母谒东君。\par
交赤松，及羡门，受要秘道爱精神。\par
食芝英，饮醴泉，柱杖桂枝佩秋兰。\par
绝人事，游浑元，若疾风游欻翩翩。\par
景未移，行数千，寿如南山不忘愆。\par

\section{秋胡行 其一}
晨上散关山，此道当何难！\par
晨上散关山，此道当何难！\par
牛顿不起，车堕谷间。\par
坐磐石之上，弹五弦之琴。\par
作为清角韵，意中迷烦。\par
歌以言志，晨上散关山。\par
有何三老公，卒来在我旁？\par
有何三老公，卒来在我旁？\par
负揜被裘，似非恒人。\par
谓卿云何困苦以自怨，徨徨所欲，来到此间？\par
歌以言志，有何三老公？\par
我居昆仑山，所谓者真人。\par
我居昆仑山，所谓者真人。\par
道深有可得。\par
名山历观，遨游八极，枕石漱流饮泉。\par
沈吟不决，遂上升天。\par
歌以言志，我居昆仑山。\par
去去不可追，长恨相牵攀。\par
去去不可追，长恨相牵攀。\par
夜夜安得寐，惆怅以自怜。\par
正而不谲，辞赋依因。\par
经传所过，西来所传。\par
歌以言志，去去不可追。\par

\section{秋胡行 其二}
愿登泰华山，神人共远游。\par
愿登泰华山，神人共远游。\par
经历昆仑山，到蓬莱。\par
飘遥八极，与神人俱。\par
思得神药，万岁为期。\par
歌以言志，愿登泰华山。\par
天地何长久！人道居之短。\par
天地何长久！人道居之短。\par
世言伯阳，殊不知老；\par
赤松王乔，亦云得道。\par
得之未闻，庶以寿考。\par
歌以言志，天地何长久！\par
明明日月光，何所不光昭！\par
明明日月光，何所不光昭！\par
二仪合圣化，贵者独人不？\par
万国率土，莫非王臣。\par
仁义为名，礼乐为荣。\par
歌以言志，明明日月关。\par
四时更逝去，昼夜以成岁。\par
四时更逝去，昼夜以成岁。\par
大人先天而天弗违。\par
不戚年往，忧世不治。\par
存亡有命，虑之为蚩。\par
歌以言志，四时更逝去。\par
戚戚欲何念！欢笑意所之。\par
戚戚欲何念！欢笑意所之。\par
壮盛智愚，殊不再来。\par
爱时进趣，将以惠谁？\par
泛泛放逸，亦同何为！\par
歌以言志，戚戚欲何念！\par

\section{善哉行 其一}
古公亶甫，积德垂仁。\par
思弘一道，哲王于豳。\par
太伯仲雍，王德之仁。\par
行施百世，断发文身。\par
伯夷叔齐，古之遗贤。\par
让国不用，饿殂首山。\par
智哉山甫，相彼宣王。\par
何用杜伯，累我圣贤。\par
齐桓之霸，赖得仲父。\par
后任竖刁，虫流出户。\par
晏子平仲，积德兼仁。\par
与世沈德，未必思命。\par
仲尼之世，主国为君。\par
随制饮酒，扬波使官。\par

\section{善哉行 其二}
自惜身薄祜，夙贱罹孤苦。\par
既无三徙教，不闻过庭语。\par
其穷如抽裂，自以思所怙。\par
虽怀一介志，是时其能与！\par
守穷者贫贱，惋叹泪如雨。\par
泣涕于悲夫，乞活安能睹？\par
我愿于天穷，琅邪倾侧左。\par
虽欲竭忠诚，欣公归其楚。\par
快人由为叹，抱情不得叙。\par
显行天教人，谁知莫不绪。\par
我愿何时随？此叹亦难处。\par
今我将何照于光曜？释衔不如雨。\par

\section{善哉行 其三}
朝日乐相乐，酣饮不知醉。\par
悲弦激新声，长笛吹清气。\par
弦歌感人肠，四坐皆欢悦。\par
寥寥高堂上，凉风入我室。\par
持满如不盈，有德者能卒。\par
君子多苦心，所愁不但一。\par
慊慊下白屋，吐握不可失。\par
众宾饱满归，主人苦不悉。\par
比翼翔云汉，罗者安所羁？\par
冲静得自然，荣华何足为！\par

\section{步出夏门行 艳}
云行雨步，超越九江之皋。\par
临观异同，心意怀犹豫，不知当复何从？\par
经过至我碣石，心惆怅我东海。\par

\section{步出夏门行 冬十月}
孟冬十月，北风徘徊，\par
天气肃清，繁霜霏霏。\par
鹍鸡晨鸣，鸿雁南飞，\par
鸷鸟潜藏，熊罴窟栖。\par
钱鎛停置，农收积场。\par
逆旅整设，以通贾商。\par
幸甚至哉！歌以咏志。\par

\section{步出夏门行 土不同}
乡土不同，河朔隆冬。\par
流澌浮漂，舟船行难。\par
锥不入地，蘴藾深奥。\par
水竭不流，冰坚可蹈。\par
士隐者贫，勇侠轻非。\par
心常叹怨，戚戚多悲。\par
幸甚至哉！歌以咏志。\par

\section{谣俗词}
瓮中无斗储，发箧无尺缯。\par
友来从我贷，不知所以应。\par

\section{董逃歌词}
德行不亏缺，变故自难常。\par
郑康成行酒，伏地气绝；\par
郭景图命尽于园桑。\par

\end{document}